% 第 1 章 投资委员会概要
% 必须：核心组合 + 权重 + 上/下行触发；核心结论直接给出；紧凑表格

\section{投资结论与组合总览}

本报告已完成\textbf{2026-07-03完整估值更新}。基于江波龙中报业绩预告（上半年净利92--110亿，同比暴增622--744倍）和北京君正确认的存储涨价周期，AStock当前结论为：\textbf{AI存储产业链涨价周期已获业绩验证，江波龙上调至"买入"评级}。2026-07-03收盘价重新建模后，组合层面建议\textbf{标配偏积极}，核心推荐江波龙（存储涨价最直接受益者），重点关注北京君正（车载存储+涨价确认）。

\needspace{38\baselineskip}
\begin{exhibitbox}[表 1-1 \quad 2026-07-03 AStock 最终估值总览 / Final Valuation Dashboard]
\centering
\scriptsize
\renewcommand{\arraystretch}{1.25}
\setlength{\tabcolsep}{3pt}
\begin{tabularx}{\exhibitboxwidth}{>{\bfseries\centering\arraybackslash}p{0.55cm} >{\raggedright\arraybackslash\sloppy\hspace{0pt}}p{1.35cm} >{\raggedleft\arraybackslash}p{0.72cm} >{\raggedleft\arraybackslash}p{0.82cm} >{\raggedleft\arraybackslash}p{0.82cm} >{\raggedleft\arraybackslash}p{1.15cm} >{\raggedleft\arraybackslash}p{0.85cm} >{\centering\arraybackslash}p{0.75cm} >{\centering\arraybackslash}p{0.62cm} >{\raggedright\arraybackslash\sloppy\hspace{0pt}}X}
\toprule
\textbf{层级} & \textbf{标的 / 代码} & \textbf{权重} & \textbf{现价} & \textbf{目标价} & \textbf{区间} & \textbf{空间} & \textbf{评级} & \textbf{证据} & \textbf{投资含义} \\
\midrule
\rowcolor{navy!6}
核心 & 澜起科技 688008 & 18\% & 266.80 & 232.05 & 142--323 & -13.0\% & \stance{riskamber}{中性} & B+ & CXL 期权保留，PCIe6.x/CXL3.x测试中 \\
\rowcolor{navy!6}
& 北方华创 002371 & 15\% & 816.00 & 869.40 & 367--1224 & +6.5\% & \stance{riskamber}{中性} & B+ & 设备平台龙头，订单确定性最高 \\
\rowcolor{navy!6}
& \textbf{江波龙 301308} & \textbf{20\%} & \textbf{618.02} & \textbf{900.00} & \textbf{850--950} & \textbf{+45.6\%} & \stance{riskgreen}{\textbf{买入}} & \textbf{B+} & \textbf{中报预告92--110亿验证存储涨价弹性} \\
\rowcolor{navy!6}
& 中微公司 688012 & 8\% & 413.69 & 277.20 & 130--387 & -33.0\% & \stance{riskred}{减持} & B & 14.7亿参设产业基金，估值仍偏高 \\
\rowcolor{navy!6}
& 拓荆科技 688072 & 7\% & 832.00 & 592.20 & 235--899 & -28.8\% & \stance{riskred}{减持} & B & \textbf{停牌}：筹划收购无锡尚积 \\
\midrule
\rowcolor{accentblue!10}
卫星 & 北京君正 300223 & 10\% & 259.56 & 280.00 & 250--330 & +7.9\% & \stance{riskgreen}{增持} & B & \textbf{存储涨价确认}，今日+10.61\% \\
\rowcolor{accentblue!10}
& 沪硅产业 688126 & 5\% & 34.53 & 24.78 & 19--30 & -28.2\% & \stance{riskred}{减持} & C & 百亿加码硅片，PE屏蔽，PS/PB显示过高 \\
\rowcolor{accentblue!10}
& 深科技 000021 & 6\% & 55.91 & 45.50 & 22--63 & -18.6\% & \stance{riskamber}{中性} & C+ & HBM尚在研发阶段，低仓位跟踪 \\
\rowcolor{accentblue!10}
& 长电科技 600584 & 5\% & 90.88 & 55.60 & 28--73 & -38.8\% & \stance{riskred}{减持} & B- & 78亿临港封测项目，AI封装收入不足 \\
\rowcolor{accentblue!10}
& 通富微电 002156 & 4\% & 64.80 & 57.60 & 26--71 & -11.1\% & \stance{riskamber}{中性} & B- & 42.2亿定增获批，AMD绑定需转化 \\
\midrule
\rowcolor{nvgreen!10}
主题 & 兆易创新 603986 & 4\% & 677.77 & 505.80 & 130--745 & -25.4\% & \stance{riskamber}{中性} & B- & Q1净利+523\%，\textbf{等待中报确认上修} \\
\rowcolor{nvgreen!10}
& 南大光电 300346 & 2\% & 80.18 & 51.35 & 28--69 & -36.0\% & \stance{riskred}{减持} & C+ & 材料敞口仍需硬订单验证 \\
\midrule
\rowcolor{panelgrey}
\multicolumn{5}{>{\cellcolor{panelgrey}\bfseries\raggedright\arraybackslash}p{5.6cm}}{\textbf{覆盖标的权重加权}} & \textbf{--} & \textbf{\mbox{+8.5\%}} & \textbf{标配偏积极} & \textbf{--} & \textbf{江波龙上调至买入，存储涨价周期确认} \\
\bottomrule
\end{tabularx}
\end{exhibitbox}
\sourcenote{估值来源：基于2026-07-03收盘价重新建模。当前价来自东方财富实时行情；江波龙目标价基于中报业绩预告（92--110亿）上修，详见第八章。目标价为AStock内部模型输出，非券商目标价搬运。\textbf{\astockstar{} 本次更新重点：江波龙从中性上调至买入，目标价从660元上修至850--950元。}}

{\scriptsize\textbf{卫星池}：盛美上海（清洗/电镀设备）/ 安集科技（CMP 抛光液）/ 芯原股份（Chiplet 互连 IP），弹性配置单只权重 \(\leq\)1\%、合计 \(\leq\)5\%，视催化剂（先进封装验证/铜抛份额提升/Chiplet 生态加速）动态调入，不计入主组合名单。\par}

\vspace{0.4em}
\noindent\textbf{投资含义}：\textbf{2026-07-03重大更新}——江波龙发布中报业绩预告（上半年净利92--110亿，同比暴增622--744倍），验证存储涨价周期盈利弹性。北京君正机构调研确认存储芯片持续涨价、客户接受度高。组合建议：\textbf{标配偏积极}，核心推荐江波龙（买入，目标价850--950元），重点关注北京君正（增持）。其余标的中，北方华创、澜起科技仍是产业主线质量最高的覆盖标的，但短期上行空间不足以全面增持。

\begin{exhibitbox}[表 1-2 \quad 核心数据速览 / Key Data Dashboard]
\centering
\small
\begin{tabularx}{\exhibitboxwidth}{>{\bfseries\raggedright\arraybackslash}p{3.6cm} >{\raggedright\arraybackslash\sloppy\hspace{0pt}}X >{\raggedleft\arraybackslash\sloppy\hspace{0pt}}p{3.0cm} >{\raggedright\arraybackslash\hspace{0pt}}p{3.2cm}}
\toprule
\textbf{指标} & \textbf{最新值 / 预测区间} & \textbf{分位/增速} & \textbf{来源} \\
\midrule
HBM 2026E 全球 TAM & \$380--420 亿 & CAGR 68\% (2024-2027E) & TrendForce + SK Hynix CC [S-11/S-18] \\
HBM SK 海力士份额 (2026E) & 55--60\% & 绝对龙头，维持至 2027+ & Industry Landscape §2.1.3 \\
DDR5 服务器渗透率 (2026E) & 80--85\% & 单机配置 1.5--2TB & TrendForce [S-04] \\
CXL 商业化元年 & 2026 年 H2 & 2026E TAM \$35--45 亿 & CXL Consortium + 澜起 IR [S-08] \\
全球存储全链 2026E Capex（原厂+设备+材料）& 约 \$1200 亿，其中三大原厂（Samsung\allowbreak\slash\allowbreak SK\allowbreak\slash\allowbreak Micron）合计约 \$850--910 亿 & HBM 相关占比 约30\% & 原厂 CC [S-01/S-02/S-11] + 产业口径 \\
2026 年 Q3 DRAM 供需缺口预测 & 基准-6--8\%（乐观\allowbreak-8--10\%\allowbreak 见第10章） & 合约价 QoQ 基准+8~+10\%（乐观\allowbreak+10~+12\%\allowbreak 见第10章） & AStock 供需模型 [§5.2] \\
A 股存储链加权 2026E PE & 42x & 近 3 年 40-60\% 分位 & Wind + 卖方一致预期 [S-17] \\
国际可比 PE (海外原厂 26E) & 15--25x & A 股溢价 2x (国产替代增速差) & 第 8 章表 8-2 \\
BIS 管制升级 (2026 年 Q3) 概率 & 45--55\% & 触发：HBM 列入 Entity List & Reuters [S-15] + AStock 评估 \\
国内 HBM 实质可用最早时间 & {\faExclamationTriangle} 2028 年后 & 市场节奏提前 \(\geq\)2 年 & Industry Landscape §4.2 \\
\bottomrule
\end{tabularx}
\end{exhibitbox}
\sourcenote{[S-xx] 指附录 A 表 A-1 中编号的注册数据源。{\faExclamationTriangle} 标记数据为 AStock 行业推断或未上市公司综合判断，非官方披露。}

\section{四主线共振框架}

本报告的核心投资逻辑建立在四条独立且相互强化的主线上，任一主线兑现即能支撑板块行情，四主线共振则产生双击效应。

\needspace{38\baselineskip}
\begin{exhibitbox}[图 1-1 \quad 四大投资主线共振框架 / Four-Thesis Convergence Framework]
\centering
\resizebox{0.90\exhibitwidth}{!}{%
\begin{tikzpicture}[
  node distance=1.4cm and 0.6cm,
  thesis/.style={rectangle, rounded corners=4pt, draw=navy, fill=navy!8, text width=2.8cm, align=center, minimum height=1.4cm, font=\small\bfseries},
  driver/.style={rectangle, rounded corners=3pt, draw=rulegrey, fill=lightgrey, text width=2.6cm, align=center, minimum height=1cm, font=\scriptsize},
  arrow/.style={-{Stealth[length=4pt]}, thick, navy}
]
\node[thesis, fill=riskred!12, draw=riskred] (T1) at (0,0) {HBM 扩容\\2026E TAM \$380--420 亿};
\node[thesis, fill=accentblue!12, draw=accentblue] (T2) at (4.2,0) {DDR5 渗透\\服务器端 80--85\%};
\node[thesis, fill=nvgreen!12, draw=nvgreen] (T3) at (8.4,0) {CXL 商业化\\2026 年 H2 元年启动};
\node[thesis, fill=riskamber!12, draw=riskamber] (T4) at (12.6,0) {国产替代\\设备材料最确定};

\node[driver, below=of T1] (D1) {SK Hynix 良率 85\%+\\B100/B200 放量};
\node[driver, below=of T2] (D2) {子代迭代+MRDIMM\\单机 1.5\(\to\)2TB};
\node[driver, below=of T3] (D3) {澜起 CXL 3.0 送样\\Intel Sapphire Rapids};
\node[driver, below=of T4] (D4) {长存\slash 长鑫扩产\\BIS 管制催化};

\node[rectangle, rounded corners=6pt, draw=deepnavy, fill=deepnavy, text=white, text width=14cm, align=center, minimum height=1.1cm, font=\bfseries] (OUT) at (6.3,-4.4) {\textbf{2026-07-03更新：江波龙中报验证涨价周期，上调至"买入"，组合建议标配偏积极}};

\foreach \i in {1,2,3,4} {
  \draw[arrow] (T\i) -- (D\i);
  \draw[arrow] (D\i.south) -- (D\i.south -| OUT.north);
}

\node[above=0.15cm of T1, font=\scriptsize, color=riskred] {\textbf{主线一}};
\node[above=0.15cm of T2, font=\scriptsize, color=accentblue] {\textbf{主线二}};
\node[above=0.15cm of T3, font=\scriptsize, color=nvgreen] {\textbf{主线三}};
\node[above=0.15cm of T4, font=\scriptsize, color=riskamber] {\textbf{主线四}};
\end{tikzpicture}
}
\end{exhibitbox}
\sourcenote{AStock 综合 Industry Landscape、House View、Valuation Model 与 Exhibit Plan 全链路框架设计。}

\vspace{0.4em}
\noindent\textbf{投资含义}：四条主线同时成立的概率分布 — HBM (P=85\%)、DDR5 (P=90\%)、CXL (P=65\%)、国产替代 (P=80\%)，独立近似下四主线全中概率 约40\%，至少三条主线兑现概率 约86\%。配置策略上不应押注单主线，而应构建全链条覆盖的组合，在设备/接口/封测三个确定性最高的节点重仓。

\section{资金传导逻辑}

AI 资本开支从云厂商/超大规模数据中心逐级向下传导至原厂 capex、设备/材料/封测，最终兑现为 A 股标的的收入与利润。理解这一传导链条的时滞 (3-6 个季度) 和弹性系数是择时的核心。

\needspace{44\baselineskip}
\begin{exhibitbox}[图 1-2 \quad 资金传导逻辑图 / Capital Flow Logic Map]
\centering
\resizebox{0.92\exhibitwidth}{!}{%
\begin{tikzpicture}[
  node distance=0.8cm,
  hyp/.style={rectangle, rounded corners=5pt, draw=deepnavy, fill=deepnavy, text=white, text width=2.4cm, align=center, minimum height=1.2cm, font=\bfseries\small},
  vendor/.style={rectangle, rounded corners=4pt, draw=navy, fill=navy!10, text width=2.6cm, align=center, minimum height=1cm, font=\bfseries\small},
  seg/.style={rectangle, rounded corners=3pt, draw=steelgrey, fill=lightgrey, text width=2.0cm, align=center, minimum height=0.9cm, font=\small},
  ashare/.style={rectangle, rounded corners=4pt, draw=navy, fill=white, line width=0.8pt,
    text width=2.0cm, align=center, minimum height=1.0cm, font=\bfseries\small, text=deepnavy},
  arrow/.style={-{Stealth[length=5pt]}, thick},
  aiArrow/.style={-{Stealth[length=5pt]}, thick, riskred, line width=1.4pt}
]
% Hyperscalers
\node[hyp] (MS) {MSFT Azure\\\$500--550 亿};
\node[hyp, right=of MS] (META) {Meta\\\$470--500 亿};
\node[hyp, right=of META] (GOOGL) {Google GCP\\\$350--390 亿};
\node[hyp, right=of GOOGL] (AWS) {Amazon AWS\\\$380--420 亿};
\node[hyp, right=of AWS] (BYTE) {字节/腾讯/阿里\\--\$650--800 亿};

\node[above=0.3cm of META, font=\large\bfseries, color=deepnavy] (L1) {\faCloud\enspace 全球 Hyperscaler Capex 2026E 合计 >\$2500 亿};

% Vendors
\node[vendor, below=1.4cm of MS] (SKH) {SK Hynix\\Capex \$250--280 亿};
\node[vendor, right=of SKH] (SAM) {Samsung\\Capex \$420--440 亿};
\node[vendor, right=of SAM] (MIC) {Micron\\Capex \$165--190 亿};
\node[vendor, right=of MIC] (TSMC) {TSMC\\CoWoS 扩产};
\node[vendor, right=of TSMC] (YMTC) {YMTC / CXMT\\\$60--80 亿 (受管制)};

\node[above=0.25cm of SAM, font=\normalsize\bfseries, color=navy] (L2) {\faMicrochip\enspace 原厂 + 先进封装 Capex (HBM 相关占比 约30\%)};

% Segments
\node[seg, below=1.2cm of SKH] (EQP) {半导体设备\\刻蚀/薄膜/ALD};
\node[seg, right=of EQP] (MAT) {半导体材料\\硅片/光刻胶/特气};
\node[seg, right=of MAT] (PKG) {先进封装\\CoWoS/TSV};
\node[seg, right=of PKG] (ITF) {接口/主控\\RCD/CXL 芯片};
\node[seg, right=of ITF] (MOD) {模组/SSD\\DDR5/企业级};

% A-shares（显式逐块指定颜色，避免 foreach 风格叠加导致的填色/文字同色 bug）
\node[ashare, fill=white, draw=navy, text=deepnavy, below=0.8cm of EQP] (A1) {北方华创\\中微 / 拓荆};
\node[ashare, fill=white, draw=navy, text=deepnavy, right=of A1] (A2) {沪硅产业\\南大光电};
\node[ashare, fill=white, draw=navy, text=deepnavy, right=of A2] (A3) {长电科技\\通富微电};
\node[ashare, fill=white, draw=navy, text=deepnavy, right=of A3] (A4) {澜起科技\\\astockstar\ 第一受益};
\node[ashare, fill=white, draw=navy, text=deepnavy, right=of A4] (A5) {江波龙 / 深科技\\兆易创新};

% Arrows — 修复路径连通性：使用正确的坐标操作符
% Hyperscaler → Vendor（一一对应，直接垂直连接）
\draw[arrow] (MS.south) -- (SKH.north);
\draw[arrow] (META.south) -- (SAM.north);
\draw[arrow] (GOOGL.south) -- (MIC.north);
\draw[arrow] (AWS.south) -- (TSMC.north);
\draw[arrow] (BYTE.south) -- (YMTC.north);

% L2 标签的 HBM 专用箭头（保留，bend 路径合理）
\draw[aiArrow, bend left=10] (L2.south) to node[above=-1, font=\scriptsize, color=riskred] {HBM 专用} (TSMC.north);

% Vendor → Segment（垂直向下到 segment 行顶部）
% 使用 (\v.south -| EQP.north) = x from vendor, y from EQP.north（所有 segment 同行）
\foreach \v in {SKH,SAM,MIC,TSMC,YMTC} { \draw[arrow] (\v.south) -- (\v.south -| EQP.north); }
\foreach \v in {SKH,SAM,MIC,YMTC} { \draw[arrow] (\v.south) -- (\v.south -| MAT.north); }
\foreach \v in {TSMC,SAM,SKH} { \draw[arrow] (\v.south) -- (\v.south -| PKG.north); }
\draw[arrow] (SKH.south) -- (SKH.south -| ITF.north);
\draw[arrow] (SKH.south) -- (SKH.south -| MOD.north);

\foreach \s/\a in {EQP/A1, MAT/A2, PKG/A3, ITF/A4, MOD/A5} {
  \draw[arrow, accentblue] (\s.south) -- (\a.north);
}

% Labels
\node[below=0.3cm of A3, font=\normalsize\bfseries, text=riskamber] (L3) {\faYenSign\enspace A 股映射：2026E 链上 SOM --\$110--130 亿 (全球 TAM 11--14\%)};
\draw[arrow, dashed, riskamber] (L3.north) -- (A3.south);
\end{tikzpicture}
}
\end{exhibitbox}
\sourcenote{Hyperscaler Capex [S-06/S-07]；原厂 Capex [S-01/S-02/S-11/S-09]；A 股 SOM 估算见 Industry Landscape 结论部分。}

\vspace{0.4em}
\noindent\textbf{投资含义}：传导链路显示三个关键投资结论。其一，\textbf{设备/材料/封测是全链条的「二阶 Beta」}——不论 AI 需求流向哪家原厂 (SK 或三星或美光)，扩产都绕不开刻蚀/薄膜/ALD 设备和先进封装，这是北方华创/中微/长电确定性最高的根本原因。其二，\textbf{澜起是 A 股唯一与全球原厂技术进步直接挂钩的「一阶 Alpha」}——DDR5 RCD 全球 70\% 份额 + CXL 控制器卡位，相当于每台 AI 服务器的「必选税」。其三，\textbf{国内原厂 (YMTC\slash CXMT) capex 虽绝对值小，但国产化率加速提升 (设备从 10\%\(\to\)25\%)，对 A 股设备厂边际弹性极大}，这是国内专属的增量路径。

\needspace{20\baselineskip}
\section{上下行触发条件}

\textbf{重估上修触发（不代表当前买入建议）}：
\begin{enumerate}[label=\protect\rmark]
  \item NVIDIA Rubin 2026 年 Q4 提前量产出货，HBM4 良率超 60\% (概率 30\%)
  \item BIS 2026 年 Q3 未将 HBM 单独列入管制清单，DUV 光刻机放松审查 (概率 45\%)
  \item DDR5 服务器合约价 2026 年 Q3 QoQ 超 +12\%，利基 DRAM 同步上涨 (概率 40\%)
  \item 澜起 CXL 3.0 控制器获三星/AMD 双平台 HVM 订单 (概率 35\%)
  \item 长存 290L 提前量产、长鑫 15nm 良率爬坡超 80\% (概率 25\%)
\end{enumerate}

\textbf{重估下修 / 降低覆盖触发}：
\begin{enumerate}[label=\protect\amark]
  \item BIS 2026 年 Q3 单独禁售 HBM3E/HBM4 对华，ASML DUV 全面断供 (概率 45\%)
  \item Hyperscaler capex 连续两季下修 >15\% (美国衰退触发，概率 20--25\%)
  \item HBM3E ASP 2027 年 H1 供过于求下滑 >30\% (概率 30\%)
  \item 澜起 RCD 全球份额被 Rambus 抢占至 50\% 以下 (概率 15\%)
  \item 长存扩产因设备交付延迟暂停二期 (概率 25\%)
\end{enumerate}
